\section{Discussion}
\label{sec:discussion}

\paragraph{What the evidence supports.}
Three claims survive the controls. First, the transition law is executable and its chunked
implementation matches the sequential reference and the full forward to numerical precision,
including at boundaries that cut through consolidation and summary blocks. Second, a
multiscale decay initialisation is not cosmetic: with the bias frozen, a single rate loses
long-range recall while the spectrum retains it. Third, a memory read needs an abstention path;
a sharper softmax alone does not make a store content-addressed.

\paragraph{Sharpness is not addressing.}
The Gap D result is the clearest lesson of the study. ZCA whitening reduced read entropy from
$1.26$ to $0.72$ nats and still left presence recall at chance. Entropy measures concentration,
not correctness: whitening amplifies every deviation direction, including the context noise
that should not be retrieved. The gated read adds two things that entropy does not capture: a
write path that refuses redundancy, and a read path that can return \emph{nothing}. The
confidence factor separates present ($0.56$) from absent ($0.03$) queries without direct
supervision, and the ablation delta rises by an order of magnitude. A memory operator without a
null hypothesis cannot express ``this is not in memory'', and no amount of temperature tuning
supplies one.

\paragraph{Negative results are part of the design.}
State expansion was intended to add capacity; at 66 bindings it made accuracy worse. Chunk
summaries were intended to add a global path; their measured benefit is within noise and an
earlier run had the opposite sign. Reporting these results changes the architecture's story:
the demonstrated recall mechanisms are the delta rule and the gated slot read, not the extra
state matrices or the summary ring. The summary ring remains valuable as a design element---it
is correct, chunk-invariant, and asymptotically cheaper than attention---but its quality
benefit is an open question, not a result.

\paragraph{Why there is no transformer language-model comparison.}
The experiment ladder's transformer is a width-matched control for synthetic recall, not a
language model. It was never trained on TinyStories. A matched LM comparison requires the same
tokenizer, corpus, token budget, optimiser, evaluation split, and roughly matched parameters or
FLOPs; that run has not been performed. Without it, statements such as ``SigmaStorm is better
or worse than a transformer'' are unsupported. This is the single most important missing
experiment, and it is cheap relative to the scale debate: a 21M transformer on 348M TinyStories
tokens fits the same two-T4 budget.

\paragraph{The long-context claim is architectural, not empirical.}
The state contract is real: 200{,}704 bytes at $T=512$ and $T=4096$ for the tiny preset, and
every operator is chunk-invariant. That is a prerequisite for streaming a 256K-token sequence
in bounded memory, but it is not the same as training or evaluating at 256K. The 21M run used
1{,}024-token windows with a differentiable horizon of $4\times256$ tokens in the reported
configuration. Long-context ability---retrieval over tens of thousands of tokens, stability of
the energy-controlled latent, and the interaction of the summary ring with the delta state---has
not been measured. The design target is 256K; the evidence stops far short of it.

\paragraph{Engineering cost.}
At equal width SigmaStorm is slower than the transformer control on CPU and has more
parameters. The operators are written in dense PyTorch; there are no fused kernels for the
chunkwise delta solve, the slot consolidation, or the summary scan. The gap in
\Cref{fig:cost} is therefore an implementation gap as much as an architectural one. A fair
efficiency comparison would use fused kernels and measure tokens per second at matched
parameter count and matched training budget. That work is not part of this report.

\paragraph{What would change the conclusion.}
Four experiments, in order of value: (i) a matched transformer LM baseline on the same corpus
and budget; (ii) a gated-read 21M run to confirm that the synthetic gain survives language
modeling; (iii) a capacity-stressed state-expansion run with enough steps for both variants to
converge; (iv) a summary ablation with multiple seeds and a longer budget. If (i) shows a large
deficit, the architecture needs a different argument than quality-per-token. If (ii) shows no
change, the slot memory is not used by the language model even when it can abstain, which is
itself an important negative result.
